%-----------------------------------------------------------------------------
% Glossar: Fachbegriffe und fremdsprachige Ausdrücke
%
% Es erscheinen nur Einträge, die im Text über \gls{}, \glspl{}, \glslink{}
% oder \glsadd{} referenziert werden.
%-----------------------------------------------------------------------------

\newglossaryentry{agile}{
  name={Agile Softwareentwicklung},
  description={Vorgehensweise, die Software in kurzen Abständen ausliefert, Anforderungen fortlaufend anpasst und das eigene Vorgehen regelmäßig reflektiert}
}

\newglossaryentry{agentic}{
  name={Agentisches Werkzeug},
  plural={Agentische Werkzeuge},
  description={KI-gestütztes Werkzeug, das in der Entwicklungsumgebung selbstständig mehrere Arbeitsschritte ausführt, also Dateien durchsucht, ändert und Programme iterativ testet}
}

\newglossaryentry{codegenerator}{
  name={Code-Generator},
  plural={Code-Generatoren},
  description={Werkzeug, das auf Basis eines großen Sprachmodells aus dem umgebenden Kontext Quellcode vorschlägt oder erzeugt}
}

\newglossaryentry{codereview}{
  name={Code Review},
  plural={Code Reviews},
  description={Systematische Prüfung von Quellcode durch andere Personen als die verfassende, um Fehler und Qualitätsmängel vor der Übernahme aufzudecken}
}

\newglossaryentry{debugging}{
  name={Debugging},
  description={Systematisches Auffinden, Eingrenzen und Beheben von Fehlern in einem Programm}
}

\newglossaryentry{devops}{
  name={DevOps},
  description={Zusammenführung von Softwareentwicklung (Development) und IT-Betrieb (Operations) mit dem Ziel, Änderungen häufig und zuverlässig auszuliefern}
}

\newglossaryentry{dod}{
  name={Definition of Done},
  description={Formale Beschreibung des Zustands, den ein Increment erreichen muss, um als fertiggestellt zu gelten; dient als gemeinsamer Qualitätsmaßstab des Teams}
}

\newglossaryentry{increment}{
  name={Increment},
  plural={Increments},
  description={In sich nutzbares Zwischenergebnis, das ein Team am Ende eines Sprints zur bisherigen Produktversion hinzufügt}
}

\newglossaryentry{konfidenzintervall}{
  name={Konfidenzintervall},
  description={Statistischer Wertebereich, der den geschätzten Kennwert mit einer festgelegten Wahrscheinlichkeit überdeckt und damit die Genauigkeit der Schätzung ausdrückt}
}

\newglossaryentry{opensource}{
  name={Open Source},
  description={Software, deren Quellcode öffentlich zugänglich ist und unter einer Lizenz steht, die Nutzung, Weitergabe und Veränderung erlaubt}
}

\newglossaryentry{overreliance}{
  name={Overreliance},
  description={Übermäßiges Vertrauen in die Ausgaben eines KI-Werkzeugs, das dazu führt, dass fehlerhafte Vorschläge ungeprüft übernommen werden}
}

\newglossaryentry{pairprogramming}{
  name={Pair Programming},
  description={Praktik, bei der zwei Entwickelnde gemeinsam an einem Arbeitsplatz denselben Quellcode schreiben und dabei fortlaufend die Rollen tauschen}
}

\newglossaryentry{repository}{
  name={Repository},
  plural={Repositories},
  description={Versioniert verwaltetes Verzeichnis, in dem der Quellcode eines Projekts samt seiner Änderungshistorie abgelegt wird}
}

\newglossaryentry{retrospektive}{
  name={Retrospektive},
  plural={Retrospektiven},
  description={Regelmäßiger Termin am Ende eines Sprints, in dem das Team das eigene Vorgehen überprüft und Verbesserungen beschließt}
}

\newglossaryentry{scrum}{
  name={Scrum},
  description={Rahmenwerk der agilen Softwareentwicklung, das die Arbeit in Sprints strukturiert und Rollen, Termine und Artefakte festlegt}
}

\newglossaryentry{sprint}{
  name={Sprint},
  plural={Sprints},
  description={Fester Zeitraum von höchstens einem Monat, in dem ein Scrum-Team ein nutzbares Increment fertigstellt}
}

\newglossaryentry{sprintreview}{
  name={Sprint Review},
  plural={Sprint Reviews},
  description={Termin am Ende eines Sprints, in dem das Increment gemeinsam mit den Beteiligten geprüft und das weitere Vorgehen abgestimmt wird}
}

\newglossaryentry{storypoint}{
  name={Story Point},
  plural={Story Points},
  description={Relative Schätzgröße für den Umfang einer Anforderung, die statt absoluter Zeitangaben den Aufwand im Vergleich zu Referenzanforderungen ausdrückt}
}

\newglossaryentry{unittest}{
  name={Unit-Test},
  plural={Unit-Tests},
  description={Automatisierter Test, der eine einzelne Programmeinheit isoliert gegen ihre erwartete Funktionsweise prüft}
}

\newglossaryentry{velocity}{
  name={Velocity},
  description={Kennzahl der je Iteration fertiggestellten Story Points; sie dient der realistischen Vorhersage künftiger Iterationen, nicht der Leistungsbewertung}
}

\newglossaryentry{xp}{
  name={Extreme Programming},
  description={Agiles Vorgehensmodell, das die Entwicklung über zwölf Vorgehensweisen wie Pair Programming, testgetriebene Entwicklung und gemeinsame Verantwortlichkeit für den Quellcode ordnet}
}
